\chapter{Statistiques robustes et queues épaisses}\label{ch:qm:robust-statistics-and-heavy-tails}

Une seule cotation erronée dans le flux, le cours de référence de l'euro enregistré à 1.5172 dollar au lieu de 1.1572 un jour de mars 2026, multiplie
par sept l'écart type des 250 derniers rendements quotidiens EUR/USD, de 0.34~\% à 2.43~\%, et déplace leur \omterm{def:qm:robust-statistics-and-heavy-tails:huber}{écart absolu médian}
de 1~\%. Un système de gestion des risques qui estime la volatilité avec l'écart type affiche désormais une devise aussi volatile qu'une petite
valeur de biotechnologie~; un système qui utilise l'\omterm{def:qm:robust-statistics-and-heavy-tails:huber}{écart absolu médian} n'a rien remarqué. Les mêmes données livrent une seconde leçon~: sur 27 ans
EUR/USD a dépassé quatre écarts types en une journée dix-sept fois, alors qu'une loi normale en attend moins d'une. Ce chapitre
est consacré aux \omterm{def:qm:estimation:estimator}{estimateurs} que les mauvaises données ne peuvent pas détourner et aux queues que les bonnes données contiennent bel et bien~: la \omterm{def:qm:robust-statistics-and-heavy-tails:influence}{fonction d'influence} et le
\omterm{def:qm:robust-statistics-and-heavy-tails:influence}{point de rupture}, la \omterm{def:qm:robust-statistics-and-heavy-tails:winsor}{winsorisation} et les rangs, les \omterm{def:qm:robust-statistics-and-heavy-tails:copula}{copules} et la dépendance de queue, les queues épaisses et la \omterm{def:qm:robust-statistics-and-heavy-tails:gev}{théorie des valeurs extrêmes} qui les extrapole,
et l'estimation de l'\omterm{def:qm:robust-statistics-and-heavy-tails:heavy}{indice de queue} sur des taux de change réels.

\section{Pertes robustes et fonction d'influence}

\begin{definition}[Fonction d'influence, point de rupture]\label{def:qm:robust-statistics-and-heavy-tails:influence}
Pour un \omterm{def:qm:estimation:estimator}{estimateur} écrit comme fonctionnelle $T(F)$ de la loi des données, la \emph{fonction d'influence}\index{fonction d'influence} en $x$ est
$\mathrm{IF}(x) = \lim_{\varepsilon \to 0}\bigl(T((1 - \varepsilon)F + \varepsilon\delta_x) - T(F)\bigr)/\varepsilon$~: l'effet sur l'estimation
d'une contamination infinitésimale au point $x$ (Hampel, 1974). Le \emph{point de rupture}\index{point de rupture} est la
plus petite fraction de l'échantillon qui, remplacée par des valeurs arbitraires, peut entraîner l'estimation aussi loin que l'on veut.
\end{definition}

La moyenne a $\mathrm{IF}(x) = x - \mu$, non bornée~; la variance a $(x - \mu)^2 - \sigma^2$, pire encore. Une seule observation placée assez loin
déplace l'une ou l'autre à sa guise.

\begin{proposition}[Points de rupture de la moyenne et de la médiane]\label{prop:qm:robust-statistics-and-heavy-tails:breakdown}
La moyenne empirique a un \omterm{def:qm:robust-statistics-and-heavy-tails:influence}{point de rupture} $1/n$, qui tend vers 0~; la médiane empirique a un \omterm{def:qm:robust-statistics-and-heavy-tails:influence}{point de rupture} $\lfloor(n + 1)/2\rfloor/n$, qui
tend vers $1/2$, la valeur maximale possible pour un \omterm{def:qm:estimation:estimator}{estimateur} de position équivariant.
\end{proposition}

\begin{proof}
Remplacer une observation par $y$ déplace la moyenne de $(y - x_1)/n$, non bornée en $y$. La médiane reste comprise entre la plus petite et la plus grande des
observations non contaminées tant qu'elles sont majoritaires~; une fois la moitié de l'échantillon remplacée, les valeurs contaminées peuvent être placées
de part et d'autre de la position médiane et l'entraîner. Aucun \omterm{def:qm:estimation:estimator}{estimateur} équivariant ne peut faire mieux~: si la moitié de l'échantillon est décalée de $t$, les données
sont tout aussi compatibles avec un échantillon qu'avec sa version translatée, et l'estimation ne peut suivre les deux.
\end{proof}

\begin{definition}[Perte de Huber, écart absolu médian]\label{def:qm:robust-statistics-and-heavy-tails:huber}
La \emph{perte de Huber}\index{perte de Huber} de seuil $c$ est $\rho_c(r) = r^2/2$ pour $|r| \le c$ et $c|r| - c^2/2$ sinon~; le \omterm{def:qm:estimation:m}{M-estimateur}
de Huber de position minimise $\sum_i\rho_c((x_i - \mu)/s)$ pour une échelle robuste $s$, si bien que sa fonction d'estimation $\psi_c(r) = \max(-c,
\min(c, r))$ est bornée (Huber, 1964). L'\emph{écart absolu médian}\index{ecart@écart absolu médian} est $\mathrm{MAD} =
1.4826\,\mathrm{med}_i|x_i - \mathrm{med}_jx_j|$, la constante qui le rend cohérent avec l'écart type d'une loi normale.
\end{definition}

L'\omterm{def:qm:estimation:estimator}{estimateur} de Huber est la moyenne au voisinage du centre et la médiane dans les queues. Avec $c = 1.345$, il conserve 95~\% de l'efficacité de la moyenne lorsque
les données sont normales ($(2\Phi(c) - 1)^2/\E[\psi_c(Z)^2]$), et sa \omterm{def:qm:robust-statistics-and-heavy-tails:influence}{fonction d'influence} est $\psi_c$, bornée par $c$ unités d'écart type. Le MAD a
un \omterm{def:qm:robust-statistics-and-heavy-tails:influence}{point de rupture} égal à un demi, mais il est inefficace et suppose la symétrie~; l'\omterm{def:qm:estimation:estimator}{estimateur} $Q_n$ de Rousseeuw et Croux (1993), un premier
quartile mis à l'échelle de toutes les distances deux à deux $|x_i - x_j|$, conserve le \omterm{def:qm:robust-statistics-and-heavy-tails:influence}{point de rupture} sans la symétrie. La \Cref{fig:qm:robust-statistics-and-heavy-tails:sensitivity}
montre ce que les \omterm{def:qm:robust-statistics-and-heavy-tails:influence}{fonctions d'influence} signifient pour une cotation erronée~: lorsque l'erreur sur le log-prix d'une journée passe de 0 à 27~\%, l'écart type sur 250 jours croît sans limite,
tandis que la MAD et $Q_n$ ne bougent qu'une fois, de 1.3~\% et 1.5~\%, puis plus du tout.

\begin{omfigure}
\begin{tikzpicture}
\begin{axis}[width=10cm,height=5cm,xlabel={erreur sur le log-prix d'une journée (\%)},ylabel={estimation quotidienne de l'échelle (\%)},
  tick label style={font=\small},label style={font=\small},xmin=0,xmax=30,ymin=0,ymax=2.6,grid=major,grid style={gray!25},
  legend style={font=\footnotesize,at={(0.03,0.97)},anchor=north west,fill=white,draw=none},table/col sep=comma]
\addplot[omThm,very thick] table[x=error,y=sd]{figdata/methods/15-robust-statistics-and-heavy-tails/sensitivity.csv};
\addplot[omDef,very thick] table[x=error,y=mad]{figdata/methods/15-robust-statistics-and-heavy-tails/sensitivity.csv};
\addplot[black,thick,dashed] table[x=error,y=qn]{figdata/methods/15-robust-statistics-and-heavy-tails/sensitivity.csv};
\addplot[only marks,mark=*,mark size=1.8pt,black] coordinates {(27.09,2.4345)};
\legend{écart type,MAD,$Q_n$}
\end{axis}
\end{tikzpicture}

\omcaption{Trois estimations de la volatilité quotidienne de l'EUR/USD à partir des 250 rendements des cours de référence jusqu'au 23 septembre 2026, lorsque le logarithme
du cours du 24 mars 2026 porte une erreur. Le point marqué correspond à la cotation à chiffres permutés 1.5172 au lieu de 1.1572 (une erreur de 27.1~\%)~: écart
type 2.43~\% au lieu de 0.34~\%, MAD 0.279~\% au lieu de 0.275~\%, $Q_n$ 0.305~\% au lieu de 0.300~\%. Données~: taux de référence de l'euro de la BCE
(source~: statistiques de la BCE).}\label{fig:qm:robust-statistics-and-heavy-tails:sensitivity}
\end{omfigure}

\section{Winsorisation, troncature et méthodes de rang}

\begin{definition}[Winsorisation, moyenne tronquée]\label{def:qm:robust-statistics-and-heavy-tails:winsor}
La \emph{winsorisation}\index{winsorisation} de niveau $p$ remplace les observations inférieures au quantile d'ordre $p$ par ce quantile et celles supérieures au
quantile d'ordre $(1 - p)$ par celui-ci. La $p$-\emph{moyenne tronquée}\index{moyenne tronquée} écarte les $\lfloor pn\rfloor$ plus petites et plus grandes
observations et calcule la moyenne du reste.
\end{definition}

Les deux ont un \omterm{def:qm:robust-statistics-and-heavy-tails:influence}{point de rupture} égal à $p$. La \omterm{def:qm:robust-statistics-and-heavy-tails:winsor}{winsorisation} est la correction courante pour une coupe transversale de signaux ou de rendements avant une régression
(chapitre~16)~: elle conserve le signe et le rang de chaque observation et plafonne son levier. C'est aussi une décision sur les données que le rapport doit
énoncer, car un \omterm{def:qm:estimation:sharpe}{ratio de Sharpe} winsorisé n'est pas le \omterm{def:qm:estimation:sharpe}{ratio de Sharpe} de la stratégie. Pour une cotation erronée, la meilleure correction se situe en amont~: filtrer
les transactions manifestement erronées par rapport à une bande robuste autour d'un centre robuste avant tout calcul de statistique.

\begin{definition}[Corrélation de rang, rho de Spearman, tau de Kendall]\label{def:qm:robust-statistics-and-heavy-tails:rank}
Une \emph{corrélation de rang}\index{corrélation de rang} mesure la dépendance à travers les seuls rangs des observations, de sorte qu'elle reste inchangée
par des transformations croissantes de l'une ou l'autre variable. La \emph{corrélation de rang de Spearman}\index{corrélation de rang de Spearman} est la corrélation
des rangs. Le \emph{tau de Kendall}\index{tau de Kendall} est la probabilité que deux couples indépendants soient concordants, moins la
probabilité qu'ils soient discordants, $\tau = \E[\operatorname{sign}((X - X')(Y - Y'))]$.
\end{definition}

\begin{definition}[Copule, coefficient de dépendance de queue]\label{def:qm:robust-statistics-and-heavy-tails:copula}
Une \emph{copule}\index{copule} est une fonction de répartition conjointe sur $[0, 1]^d$ à marges uniformes. Le \emph{coefficient de dépendance
de queue}\index{coefficient de dépendance de queue} inférieur de $(X, Y)$ à marges continues $F_X, F_Y$ est $\lambda_L = \lim_{q \downarrow
0}\P(F_Y(Y) \le q \mid F_X(X) \le q)$, et le coefficient supérieur se définit symétriquement.
\end{definition}

\begin{theorem}[Sklar]\label{thm:qm:robust-statistics-and-heavy-tails:sklar}
Toute fonction de répartition conjointe $H$ de marges $F_1, \dots, F_d$ peut s'écrire $H(x_1, \dots, x_d) = C(F_1(x_1), \dots, F_d(x_d))$ pour une
\omterm{def:qm:robust-statistics-and-heavy-tails:copula}{copule} $C$, unique lorsque les marges sont continues~; réciproquement, toute \omterm{def:qm:robust-statistics-and-heavy-tails:copula}{copule} et toutes les marges définissent ainsi une loi conjointe.
\end{theorem}

Pour des marges continues, la preuve tient en une ligne~: $C$ est la loi de $(F_1(X_1), \dots, F_d(X_d))$, dont les marges sont uniformes. Le théorème
sépare ce qu'une loi conjointe dit de chaque variable de ce qu'elle dit de leur dépendance, et les \omterm{def:qm:robust-statistics-and-heavy-tails:rank}{corrélations de rang} et la dépendance de queue
sont des propriétés de la seule \omterm{def:qm:robust-statistics-and-heavy-tails:copula}{copule}. Pour la \omterm{def:qm:robust-statistics-and-heavy-tails:copula}{copule} gaussienne de corrélation $\rho$, le rho de Spearman vaut $\frac6\pi\arcsin\frac\rho2$, le
\omterm{def:qm:robust-statistics-and-heavy-tails:rank}{tau de Kendall} vaut $\frac2\pi\arcsin\rho$, et les deux \omterm{def:qm:robust-statistics-and-heavy-tails:copula}{coefficients de dépendance de queue} sont nuls~: les chutes conjointes deviennent asymptotiquement indépendantes, quelle que soit
la valeur de $\rho$. Les modèles de portefeuille de crédit héritent de cette propriété de la \omterm{def:qm:robust-statistics-and-heavy-tails:copula}{copule} gaussienne, un point que le livre~6 reprend avec la \omterm{def:qm:robust-statistics-and-heavy-tails:copula}{copule} de
Student-$t$, dont la dépendance de queue est positive.

Le dollar et la livre sterling, tous deux cotés en euros, présentent une corrélation de Pearson quotidienne de 0.42 sur 1999--2026. Le rho de Spearman vaut 0.406, exactement
ce qu'implique une \omterm{def:qm:robust-statistics-and-heavy-tails:copula}{copule} gaussienne de même corrélation (0.406), et un \omterm{def:qm:robust-statistics-and-heavy-tails:rank}{tau de Kendall} de 0.286, contre 0.277 implicite. Les queues divergent~: sur les
5~\% de jours où le dollar a le plus reculé face à l'euro, la livre sterling se trouvait parmi ses propres 5~\% de pires jours dans 32~\% des cas, contre 20~\% pour la
\omterm{def:qm:robust-statistics-and-heavy-tails:copula}{copule} gaussienne~; à 1~\%, 21~\% contre 10~\% (\cref{fig:qm:robust-statistics-and-heavy-tails:copula}). La dépendance moyenne est gaussienne~; les
extrêmes conjoints ne le sont pas.

\begin{omfigure}
\begin{tikzpicture}
\begin{axis}[width=7cm,height=7cm,xlabel={rang du rendement USD par EUR},ylabel={rang du rendement GBP par EUR},
  tick label style={font=\small},label style={font=\small},xmin=0,xmax=1,ymin=0,ymax=1,
  table/col sep=comma]
\addplot[only marks,mark=*,mark size=0.35pt,omDef!70] table[x=u,y=v]{figdata/methods/15-robust-statistics-and-heavy-tails/copula.csv};
\draw[omThm,thick] (axis cs:0,0) rectangle (axis cs:0.05,0.05);
\draw[omThm,thick] (axis cs:0.95,0.95) rectangle (axis cs:1,1);
\end{axis}
\end{tikzpicture}

\omcaption{La \omterm{def:qm:robust-statistics-and-heavy-tails:copula}{copule} empirique des rendements quotidiens de l'euro face au dollar et à la livre sterling, 1999--2026 (un jour sur quatre représenté)~: chaque point
est le couple des rangs du jour, ramené à $(0, 1)$. Les carrés marquent les queues conjointes à 5~\%, où 32~\% des pires jours du dollar figurent aussi
parmi les pires de la livre sterling, contre 20~\% pour une \omterm{def:qm:robust-statistics-and-heavy-tails:copula}{copule} gaussienne de même corrélation. Données~: taux de référence de l'euro de la BCE (source~:
statistiques de la BCE).}\label{fig:qm:robust-statistics-and-heavy-tails:copula}
\end{omfigure}

\section{Queues épaisses}

\begin{definition}[Loi à queue épaisse, indice de queue]\label{def:qm:robust-statistics-and-heavy-tails:heavy}
Une \emph{loi à queue épaisse}\index{loi à queue épaisse} (au sens de Pareto) vérifie $\P(X > x) = x^{-\alpha}L(x)$ avec $L$ à
variation lente ($L(tx)/L(x) \to 1$ lorsque $x \to \infty$, pour tout $t > 0$). L'exposant $\alpha > 0$ est son \emph{indice de queue}\index{indice de queue}.
\end{definition}

\begin{proposition}[Moments au-delà de l'indice de queue]\label{prop:qm:robust-statistics-and-heavy-tails:moments}
Si $X \ge 0$ a pour \omterm{def:qm:robust-statistics-and-heavy-tails:heavy}{indice de queue} $\alpha$, alors $\E[X^\beta] < \infty$ pour $\beta < \alpha$ et $\E[X^\beta] = \infty$ pour $\beta > \alpha$.
\end{proposition}

\begin{proof}
$\E[X^\beta] = \int_0^\infty\beta x^{\beta-1}\P(X > x)\,dx$. Pour $x$ grand, une fonction $L$ à variation lente vérifie $x^{-\varepsilon} \le L(x) \le
x^\varepsilon$ pour tout $\varepsilon > 0$ (bornes de Potter), de sorte que l'intégrande est comprise entre $x^{\beta - \alpha - 1 \mp \varepsilon}$~: intégrable
à l'infini lorsque $\beta < \alpha$ (en prenant $\varepsilon$ petit) et non lorsque $\beta > \alpha$.
\end{proof}

Une loi de Student $t$ à $\nu$ degrés de liberté a pour \omterm{def:qm:robust-statistics-and-heavy-tails:heavy}{indice de queue} $\nu$. L'étude de Mandelbrot de 1963 sur les prix du coton proposait des lois stables, dont la
variance est infinie~; le graphe de Hill ci-dessous donne aux pertes quotidiennes de EUR/USD un \omterm{def:qm:robust-statistics-and-heavy-tails:heavy}{indice de queue} voisin de quatre, d'où une variance finie et un
quatrième moment peu fiable, conformément à l'étude de Cont sur les faits stylisés, qui trouve pour la plupart des séries de rendements étudiées des \omterm{def:qm:robust-statistics-and-heavy-tails:heavy}{indices de queue} supérieurs à deux et inférieurs à cinq. La conséquence pratique est que le kurtosis, moment d'ordre quatre, est estimé mal, voire pas du tout~: celui de EUR/USD
sur 27 ans vaut 6.3, et quelques jours suffisent à le fixer. Un graphe quantile-quantile montre directement les queues
(\cref{fig:qm:robust-statistics-and-heavy-tails:qq}).

\begin{omfigure}
\begin{tikzpicture}
\begin{axis}[width=8cm,height=6.2cm,xlabel={quantile normal},ylabel={rendement EUR/USD standardisé},
  tick label style={font=\small},label style={font=\small},xmin=-4.2,xmax=4.2,ymin=-8.5,ymax=8.5,grid=major,grid style={gray!25},
  legend style={font=\footnotesize,at={(0.03,0.97)},anchor=north west,fill=white,draw=none},table/col sep=comma]
\addplot[only marks,mark=*,mark size=0.9pt,omDef] table[x=normal,y=sample]{figdata/methods/15-robust-statistics-and-heavy-tails/qq.csv};
\addplot[black,dashed,thick,domain=-4.2:4.2,samples=2] {x};
\legend{rendements quotidiens 1999--2026,loi normale}
\end{axis}
\end{tikzpicture}

\omcaption{Graphe quantile-quantile des 7\,098 rendements quotidiens de EUR/USD, standardisés, face à la loi normale (chaque observation dans les
queues, une sur cinquante dans le corps de la loi). Les plus fortes variations, $-8.2$ et $+7.2$ écarts-types, se situent là où une loi normale attribue une cote
inférieure à un sur cent milliards. Données~: taux de référence de l'euro de la BCE (source~: statistiques de la BCE).}\label{fig:qm:robust-statistics-and-heavy-tails:qq}
\end{omfigure}

\section{Théorie des valeurs extrêmes}

\begin{definition}[Théorie des valeurs extrêmes, loi GEV]\label{def:qm:robust-statistics-and-heavy-tails:gev}
La \emph{théorie des valeurs extrêmes}\index{théorie des valeurs extrêmes} étudie les lois des maxima et des dépassements de seuils élevés. La
\emph{loi des valeurs extrêmes généralisée}\index{loi des valeurs extrêmes généralisée} de paramètre de forme $\xi$ a pour fonction de répartition
$G_\xi(x) = \exp\bigl(-(1 + \xi x)^{-1/\xi}\bigr)$ pour $1 + \xi x > 0$ (et $\exp(-e^{-x})$ pour $\xi = 0$), à un paramètre de position et d'échelle près.
\end{definition}

\begin{theorem}[Fisher--Tippett--Gnedenko]\label{thm:qm:robust-statistics-and-heavy-tails:ftg}
Si, pour des $X_i$ i.i.d. et des constantes $a_n > 0$, $b_n$, la loi de $(\max_{i \le n}X_i - b_n)/a_n$ converge vers une limite non dégénérée, cette limite est une
\omterm{def:qm:robust-statistics-and-heavy-tails:gev}{loi des valeurs extrêmes généralisée}. Des queues épaisses d'indice $\alpha$ donnent $\xi = 1/\alpha > 0$~; des queues normales, exponentielles et similaires
donnent $\xi = 0$~; des queues bornées donnent $\xi < 0$.
\end{theorem}

\begin{definition}[Loi de Pareto généralisée, méthode des excès au-delà d'un seuil]\label{def:qm:robust-statistics-and-heavy-tails:gpd}
La \emph{loi de Pareto généralisée}\index{loi de Pareto généralisée} de paramètre de forme $\xi$ et d'échelle $\beta$ vérifie $\P(Y > y) = (1 +
\xi y/\beta)^{-1/\xi}$ pour $y \ge 0$ (et $e^{-y/\beta}$ pour $\xi = 0$). La \emph{méthode des excès au-delà d'un seuil}\index{méthode des excès au-delà d'un seuil} (POT) ajuste
par maximum de vraisemblance la loi aux excès $X - u$ des observations au-delà d'un seuil élevé $u$.
\end{definition}

\begin{theorem}[Pickands--Balkema--de Haan]\label{thm:qm:robust-statistics-and-heavy-tails:pbdh}
Pour une loi appartenant au domaine d'attraction de $G_\xi$, la loi de l'excès $X - u$ sachant $X > u$ est approchée, uniformément en
l'excès, par une \omterm{def:qm:robust-statistics-and-heavy-tails:gpd}{loi de Pareto généralisée} de forme $\xi$ et d'échelle $\beta(u)$ quelconque, l'erreur tendant vers zéro lorsque $u$ tend vers la borne
supérieure du support.
\end{theorem}

Ce théorème autorise une extrapolation~: avec $N_u$ des $n$ observations au-dessus de $u$, la perte dépassée avec probabilité $1 - p$ est
\[
\mathrm{VaR}_p = u + \frac\beta\xi\Bigl(\bigl(\tfrac{N_u/n}{1 - p}\bigr)^\xi - 1\Bigr).
\]
Pour les pertes quotidiennes de EUR/USD au-dessus de leur quantile à 95~\% ($u = 0.93\%$, 355 excès), le maximum de vraisemblance donne $\xi = 0.069$ et $\beta = 0.34\%$,
et la perte d'un jour sur mille est de 2.47~\%. Le $\xi$ ajusté dépend du seuil (0.03 au-dessus du quantile à 90~\%, 0.21 au-dessus du
99~\%), le quantile beaucoup moins (2.43~\% à 2.49~\% pour les mêmes seuils)~: l'extrapolation est plus stable que le paramètre.

\section{Estimation de l'indice de queue}

\begin{definition}[Estimateur de Hill]\label{def:qm:robust-statistics-and-heavy-tails:hill}
À partir des statistiques d'ordre $X_{(1)} \ge X_{(2)} \ge \dots$ d'un échantillon positif, l'\emph{estimateur de Hill}\index{estimateur de Hill} de l'\omterm{def:qm:robust-statistics-and-heavy-tails:heavy}{indice de queue}
issu des $k$ plus grandes valeurs est $\hat\alpha_k = \bigl(\frac1k\sum_{i=1}^k\ln X_{(i)} - \ln X_{(k+1)}\bigr)^{-1}$ (Hill, 1975).
\end{definition}

C'est l'\omterm{def:qm:estimation:mle}{estimateur du maximum de vraisemblance} de $\alpha$ pour une queue de Pareto exacte au-delà de $X_{(k+1)}$, et, lorsque $k \to \infty$ avec $k/n \to 0$ à un
rythme convenable, $\sqrt k(\hat\alpha_k - \alpha) \to \mathcal N(0, \alpha^2)$. Le choix de $k$ relève d'un compromis biais-variance~: un petit $k$ est bruité,
un grand $k$ empiète sur le corps de la loi, là où la forme de Pareto cesse de tenir. La courbe de Hill illustre ce compromis
(\cref{fig:qm:robust-statistics-and-heavy-tails:hill})~: pour les pertes de EUR/USD, $\hat\alpha$ vaut 4.9 pour $k = 50$ ($\pm 0.7$), proche de 3.9 pour $k =
100$ à 200, puis descend jusqu'à 2.6 pour $k = 700$, où l'on ajuste le corps de la loi.

\begin{omfigure}
\begin{tikzpicture}
\begin{axis}[width=10cm,height=5cm,xlabel={nombre des plus grandes pertes utilisées, $k$},ylabel={Hill $\hat\alpha_k$},
  tick label style={font=\small},label style={font=\small},xmin=20,xmax=700,ymin=0,ymax=8,grid=major,grid style={gray!25},
  legend style={font=\footnotesize,at={(0.97,0.97)},anchor=north east,fill=white,draw=none},table/col sep=comma]
\addplot[name path=hi,draw=none,forget plot] table[x=k,y=hi]{figdata/methods/15-robust-statistics-and-heavy-tails/hill.csv};
\addplot[name path=lo,draw=none,forget plot] table[x=k,y=lo]{figdata/methods/15-robust-statistics-and-heavy-tails/hill.csv};
\addplot[omDef!20,forget plot] fill between[of=hi and lo];
\addplot[omDef,very thick] table[x=k,y=alpha]{figdata/methods/15-robust-statistics-and-heavy-tails/hill.csv};
\legend{estimation de Hill (bande grisée à $\pm 2$ erreurs types)}
\end{axis}
\end{tikzpicture}

\omcaption{Graphe de Hill des pertes quotidiennes de EUR/USD, 1999--2026~: l'\omterm{def:qm:robust-statistics-and-heavy-tails:heavy}{indice de queue} estimé à partir des $k$ plus grandes pertes, avec deux
\omterm{def:qm:estimation:estimator}{erreurs types} asymptotiques. Un palier autour de 3.9 pour $k$ entre 100 et 200~; au-delà, l'estimation dérive à mesure que le corps de la distribution entre en jeu.
Données~: taux de référence de l'euro de la BCE (source~: statistiques de la BCE).}\label{fig:qm:robust-statistics-and-heavy-tails:hill}
\end{omfigure}

Trois modèles répondent alors à la question du gestionnaire des risques~: la perte qui ne se produit qu'une fois sur mille jours (la \cref{fig:qm:robust-statistics-and-heavy-tails:tail}).
La loi normale, avec la volatilité échantillonnale de 0.58~\%, la place à 1.80~\%, et 35 jours sur les 7\,098 ont subi une perte supérieure, contre 7 attendus. Une loi de Student $t$
ajustée par maximum de vraisemblance ($\nu = 4.7$, échelle 0.44~\%) la place à 2.74~\%. L'ajustement par \omterm{def:qm:robust-statistics-and-heavy-tails:gpd}{loi de Pareto généralisée} la place à 2.47~\%, dépassée lors de 6 jours~;
le quantile empirique vaut 2.27~\%.

\begin{omfigure}
\begin{tikzpicture}
\begin{semilogyaxis}[width=10cm,height=5.4cm,xlabel={perte quotidienne de EUR/USD (\%)},ylabel={probabilité d'une perte plus grande},
  tick label style={font=\small},label style={font=\small},xmin=0.5,xmax=4.5,ymin=1e-5,ymax=0.5,grid=major,grid style={gray!25},
  unbounded coords=jump,
  legend columns=4,legend style={font=\footnotesize,at={(0.5,-0.3)},anchor=north,draw=none,fill=none,
  /tikz/every even column/.append style={column sep=6pt}},table/col sep=comma]
\addplot[only marks,mark=*,mark size=1.3pt,black] table[x=loss,y=empirical]{figdata/methods/15-robust-statistics-and-heavy-tails/tail.csv};
\addplot[omThm,very thick,dashed] table[x=loss,y=normal]{figdata/methods/15-robust-statistics-and-heavy-tails/tail.csv};
\addplot[black!60,very thick,dotted] table[x=loss,y=student]{figdata/methods/15-robust-statistics-and-heavy-tails/tail.csv};
\addplot[omDef,very thick] table[x=loss,y=gpd]{figdata/methods/15-robust-statistics-and-heavy-tails/tail.csv};
\addplot[black!40,thin,sharp plot] coordinates {(0.5,0.001) (4.5,0.001)};
\legend{données,normale,Student $t$,{GPD au-dessus de 0.93~\%}}
\end{semilogyaxis}
\end{tikzpicture}

\omcaption{Probabilités de dépassement des pertes quotidiennes de EUR/USD, 1999--2026, face à trois ajustements~: la loi normale, une loi de Student $t$ ($\nu = 4.7$) et une
\omterm{def:qm:robust-statistics-and-heavy-tails:gpd}{loi de Pareto généralisée} au-dessus du quantile à 95~\%. La ligne grise est le niveau d'une perte sur mille jours, franchi à 1.80~\%, 2.74~\% et 2.47~\%. Données~:
taux de référence de l'euro de la BCE (source~: statistiques de la BCE).}\label{fig:qm:robust-statistics-and-heavy-tails:tail}
\end{omfigure}

Intégrons maintenant la transaction erronée de l'accroche dans l'historique complet. Le quantile normal monte à 2.27~\% (+26~\%), car une paire de rendements de $\pm 27\%$
gonfle l'écart-type d'un quart. Le quantile de Pareto généralisé monte à 2.85~\% (+15~\%)~: la mauvaise perte est de loin le plus grand excès
et fait passer le paramètre de forme ajusté de 0.07 à 0.22. La loi de Student $t$ varie le moins, à 2.91~\% (+6~\%), sa vraisemblance bornant l'influence de
tout point isolé. Aucune estimation de queue n'est robuste au sens de la première section~; la queue est l'endroit où l'information est rare, et une
seule observation en constitue une large part. La parade est en amont, dans le filtre.

\section{Tutoriel~: le mauvais tick}\label{tut:qm:robust-statistics-and-heavy-tails:badtick}

\begin{tutorial}
\textbf{Objectif.} Estimer de façon robuste la volatilité et la queue de EUR/USD, puis altérer une transaction et mesurer ce qui varie.
\textbf{État final~:} \cref{fig:qm:robust-statistics-and-heavy-tails:sensitivity,fig:qm:robust-statistics-and-heavy-tails:hill,fig:qm:robust-statistics-and-heavy-tails:tail}
et les trois pertes qui ne surviennent qu'une fois sur mille jours, avec et sans la mauvaise transaction.
\begin{tutsteps}
\item \textbf{L'\omterm{def:qm:robust-statistics-and-heavy-tails:hill}{estimateur de Hill}} à partir des $k$ plus grandes pertes.
\omcode{code/firm/robust/firm_robust.py}{124}{132}{L'estimateur de Hill avec son erreur type asymptotique.}\label{lst:qm:robust-statistics-and-heavy-tails:hill}
\item \textbf{La vraisemblance de la \omterm{def:qm:robust-statistics-and-heavy-tails:gpd}{loi de Pareto généralisée}, l'ajustement et le quantile.}
\omcode{code/firm/robust/firm_robust.py}{166}{192}{Méthode des excès au-delà d'un seuil (POT)~: vraisemblance, ajustement et quantile.}\label{lst:qm:robust-statistics-and-heavy-tails:gpd}
\item \textbf{Exécuter} \texttt{scale\_table(False)}, \texttt{scale\_table(True)}, \texttt{tail\_quantiles(False)}, \texttt{tail\_quantiles(True)} et
\texttt{dependence()} dans \texttt{qm\_robust.py}, puis \texttt{fig\_robust.py}.
\end{tutsteps}
\textbf{Pour aller plus loin.} Construire le filtre~: signaler une transaction dont le rendement dépasse dix fois l'\omterm{def:qm:robust-statistics-and-heavy-tails:huber}{écart absolu médian} (MAD) des 60 derniers jours et qui s'inverse le
lendemain, et vérifier qu'il détecte la mauvaise transaction sans signaler 2008~; estimer la queue supérieure (gains) et comparer.
\end{tutorial}

\section{Construction~: le module de statistiques robustes}\label{bld:qm:robust-statistics-and-heavy-tails:robust}

\begin{build}
\bfield{Rôle} Échelles robustes pour toute volatilité que la firme miniature estime à partir des transactions, et les estimations de queue sur lesquelles reposent ses limites de risque.
\bfield{Interface} \texttt{mad}, \texttt{qn\_scale}, \texttt{huber\_location}~; \texttt{winsorize}, \texttt{trimmed\_mean}~; \texttt{spearman},
\texttt{kendall\_tau}, \texttt{tail\_dependence(x, y, q)}~; \texttt{hill(x, k)}~; \texttt{gpd\_fit(excess)}, \texttt{gpd\_quantile(u, xi, beta,
p\_u, p)}.
\bfield{Règles} Les échelles sont rapportées avec l'\omterm{def:qm:estimation:estimator}{estimateur} nommé~; un quantile de queue est rapporté avec son seuil, son nombre d'excès et sa
sensibilité au seuil~; aucune mémoire en $O(n^2)$ pour les statistiques par paires.
\bfield{Tests d'acceptation} \texttt{code/firm/robust/tests/}~: MAD et $Q_n$ convergents pour la loi normale et stables sous 10~\% d'erreurs grossières~; $Q_n$ égal
à sa définition par force brute~; les formules de \omterm{def:qm:robust-statistics-and-heavy-tails:rank}{corrélation de rang} de la \omterm{def:qm:robust-statistics-and-heavy-tails:copula}{copule} gaussienne~; Hill sur une loi de Pareto exacte~; l'ajustement GPD retrouve ses
paramètres et son quantile est dépassé au taux annoncé.
\bfield{Pour aller plus loin} Maxima par blocs avec la \omterm{def:qm:robust-statistics-and-heavy-tails:gev}{loi des valeurs extrêmes généralisée}~; une \omterm{def:qm:robust-statistics-and-heavy-tails:gpd}{méthode des excès au-delà d'un seuil} avec déclustérisation pour les pertes dépendantes~; la \omterm{def:qm:robust-statistics-and-heavy-tails:copula}{copule} de Student-$t$ ajustée au
dollar et à la livre sterling.
\end{build}

\begin{omsources}
\item R.~Cont, «~Empirical properties of asset returns: stylized facts and statistical issues~», \emph{Quantitative Finance} 1, 2001.
\item P.~J.~Huber, «~Robust estimation of a location parameter~», \emph{Annals of Mathematical Statistics} 35, 1964.
\item F.~R.~Hampel, «~The influence curve and its role in robust estimation~», \emph{Journal of the American Statistical Association} 69, 1974.
\item P.~J.~Rousseeuw et C.~Croux, «~Alternatives to the median absolute deviation~», \emph{Journal of the American Statistical Association}
88, 1993.
\item A.~Sklar, «~Fonctions de répartition à $n$ dimensions et leurs marges~», \emph{Publications de l'Institut de Statistique de
l'Université de Paris} 8, 1959.
\item R.~A.~Fisher et L.~H.~C.~Tippett, «~Limiting forms of the frequency distribution of the largest or smallest member of a sample~»,
\emph{Proceedings of the Cambridge Philosophical Society} 24, 1928~; B.~Gnedenko, \emph{Annals of Mathematics} 44, 1943.
\item A.~A.~Balkema et L.~de Haan, «~Residual life time at great age~», \emph{Annals of Probability} 2, 1974~; J.~Pickands, «~Statistical
inference using extreme order statistics~», \emph{Annals of Statistics} 3, 1975.
\item B.~M.~Hill, «~A simple general approach to inference about the tail of a distribution~», \emph{Annals of Statistics} 3, 1975.
\item B.~Mandelbrot, «~The variation of certain speculative prices~», \emph{Journal of Business} 36, 1963.
\item European Central Bank, euro foreign exchange reference rates against the dollar and the pound (ECB Data Portal, série
\texttt{EXR.D.USD.EUR.SP00.A} et \texttt{EXR.D.GBP.EUR.SP00.A}), consulté le 24 septembre 2026.
\end{omsources}

\section{Exercices}

\begin{exercise}[$\star$]\label{exo:qm:robust-statistics-and-heavy-tails:1}
Quels sont les \omterm{def:qm:robust-statistics-and-heavy-tails:influence}{points de rupture} de la \omterm{def:qm:robust-statistics-and-heavy-tails:winsor}{moyenne tronquée} à 10~\% et de l'écart interquartile~?
\end{exercise}

\begin{exercise}[$\star$]\label{exo:qm:robust-statistics-and-heavy-tails:2}
Une queue de Pareto a un \omterm{def:qm:robust-statistics-and-heavy-tails:heavy}{indice de queue} égal à 3. Parmi la moyenne, la variance et le kurtosis, lesquels existent~?
\end{exercise}

\begin{exercise}[$\star$]\label{exo:qm:robust-statistics-and-heavy-tails:3}
Pour une \omterm{def:qm:robust-statistics-and-heavy-tails:copula}{copule} gaussienne de corrélation 0.42, calculer le rho de Spearman et le \omterm{def:qm:robust-statistics-and-heavy-tails:rank}{tau de Kendall}.
\end{exercise}

\begin{exercise}[$\star\star$]\label{exo:qm:robust-statistics-and-heavy-tails:4}
Montrer que la fonction d'estimation de Huber avec $c = 1.345$ a une efficacité de 95~\% pour la loi normale~: évaluer $(2\Phi(c) - 1)^2/\E[\psi_c(Z)^2]$.
\end{exercise}

\begin{exercise}[$\star\star$]\label{exo:qm:robust-statistics-and-heavy-tails:5}
Montrer que l'\omterm{def:qm:robust-statistics-and-heavy-tails:hill}{estimateur de Hill} est l'\omterm{def:qm:estimation:mle}{estimateur du maximum de vraisemblance} de $\alpha$ pour les $k$ plus grandes observations d'une loi de Pareto exacte au-dessus de
$X_{(k+1)}$.
\end{exercise}

\begin{exercise}[$\star\star$]\label{exo:qm:robust-statistics-and-heavy-tails:6}
Avec $u = 0.93\%$, 355 excès sur 7\,098 jours, $\xi = 0.069$ et $\beta = 0.34\%$, calculer la perte dépassée une fois tous les dix mille jours.
\end{exercise}

\begin{exercise}[$\star\star\star$]\label{exo:qm:robust-statistics-and-heavy-tails:7}
\emph{Programmation.} Réajuster la \omterm{def:qm:robust-statistics-and-heavy-tails:gpd}{loi de Pareto généralisée} aux pertes EUR/USD au-delà des quantiles à 90~\%, 95~\%, 97.5~\% et 99~\%, et présenter dans un tableau $\xi$ et la
perte dépassée une fois tous les mille jours.
\end{exercise}

\begin{exercise}[$\star\star\star$]\label{exo:qm:robust-statistics-and-heavy-tails:8}
\emph{Repérer la faille.} «~Notre modèle de risque estime la volatilité de chaque devise avec l'écart-type des 250 derniers jours, mais nous
winsorisons d'abord les rendements aux quantiles à 1~\% et 99~\%, afin que les mauvaises cotations ne puissent pas nous nuire.»
\end{exercise}

\section{Problème~: la mauvaise cotation}

\begin{problem}[{Problème du week-end --- une cotation erronée et la perte d'une fois sur mille jours}]\label{pb:qm:robust-statistics-and-heavy-tails:1}
Les données sont les taux de référence quotidiens de l'euro contre le dollar publiés par la BCE, du 4 janvier 1999 au 23 septembre 2026 (7\,098 rendements quotidiens). Dans la
version corrompue, la cotation du 24 mars 2026 affiche 1.5172 au lieu de 1.1572.

\medskip\noindent\textbf{Partie I --- Ordre de grandeur.}
\begin{enumerate}
\item Quelle erreur de log-prix la mauvaise cotation introduit-elle, et quels rendements touche-t-elle~?
\item Quels sont l'écart-type, l'\omterm{def:qm:robust-statistics-and-heavy-tails:huber}{écart absolu médian} (MAD) et $Q_n$ des 250 derniers rendements, propres et corrompus~?
\item Quelles sont la moyenne, la médiane et l'estimation de position de Huber de ces rendements, propres et corrompus~?
\item Combien de jours de l'historique complet ont enregistré une perte supérieure à quatre écarts-types, à comparer à l'espérance d'une loi normale~?
\item Quel est le coefficient d'aplatissement empirique, et pourquoi est-ce un mauvais résumé~?
\end{enumerate}

\medskip\noindent\textbf{Partie II --- La queue.}
\begin{enumerate}[resume]
\item Que donne le graphe de Hill pour l'\omterm{def:qm:robust-statistics-and-heavy-tails:heavy}{indice de queue}, et pour quelles valeurs de seuil~?
\item Quels sont le seuil, le nombre d'excès et les $\xi$ et $\beta$ ajustés par la \omterm{def:qm:robust-statistics-and-heavy-tails:gpd}{méthode des excès au-delà d'un seuil} (POT)~?
\item Quels degrés de liberté et quelle échelle donne l'ajustement par une loi de Student-$t$~?
\item Quelles sont les trois pertes d'une fois sur mille jours, et la perte empirique~?
\item Combien de jours ont dépassé les quantiles normal et GPD, à comparer aux 7 attendus~?
\end{enumerate}

\medskip\noindent\textbf{Partie III --- La mauvaise cotation dans la queue.}
\begin{enumerate}[resume]
\item Que deviennent les trois quantiles avec la mauvaise cotation~?
\item Pourquoi le quantile normal se déplace-t-il le plus~?
\item Que devient le $\xi$ ajusté~?
\item Pourquoi la loi de Student $t$ se déplace-t-elle le moins~?
\item Une échelle robuste dans le modèle normal l'aurait-elle sauvé~?
\end{enumerate}

\medskip\noindent\textbf{Partie IV --- Jugement.}
\begin{enumerate}[resume]
\item Quel quantile faudrait-il retenir dans le système de risque, et sous quelle réserve~?
\item Où la mauvaise cotation aurait-elle dû être arrêtée~?
\item Comment la dépendance de queue entre le dollar et la livre modifie-t-elle une limite portant sur deux devises~?
\item Énoncer le \emph{résultat nommé}~: les trois quantiles et l'erreur qu'une seule mauvaise cotation introduit dans chacun.
\item En une phrase~: qu'est-ce qui peut être rendu robuste, et qu'est-ce qui ne le peut pas~?
\end{enumerate}
\end{problem}

\section{Questions d'entretien}

\begin{interviewq}[$\star$ \iqroles{chercheur, développeur}]\label{iq:qm:robust-statistics-and-heavy-tails:1}
Votre estimation de volatilité d'une devise calme est multipliée par sept en une nuit. Que vérifiez-vous en premier~?
\end{interviewq}

\begin{interviewq}[$\star\star$ \iqroles{chercheur, risque}]\label{iq:qm:robust-statistics-and-heavy-tails:2}
Qu'est-ce qu'un \omterm{def:qm:robust-statistics-and-heavy-tails:influence}{point de rupture}~? Donnez les \omterm{def:qm:robust-statistics-and-heavy-tails:influence}{points de rupture} de la moyenne, de la médiane et du MAD.
\end{interviewq}

\begin{interviewq}[$\star\star$ \iqroles{chercheur, risque}]\label{iq:qm:robust-statistics-and-heavy-tails:3}
Qu'implique un \omterm{def:qm:robust-statistics-and-heavy-tails:heavy}{indice de queue} égal à 3 pour les moments des rendements, et pour l'estimation du coefficient d'aplatissement~?
\end{interviewq}

\begin{interviewq}[$\star\star$ \iqroles{risque, chercheur}]\label{iq:qm:robust-statistics-and-heavy-tails:4}
Comment estimeriez-vous le quantile à 99.9~\% de la perte quotidienne à partir de 25 ans de données~?
\end{interviewq}

\begin{interviewq}[$\star\star$ \iqroles{chercheur, ingénieur ML}]\label{iq:qm:robust-statistics-and-heavy-tails:5}
Pourquoi utiliser une \omterm{def:qm:robust-statistics-and-heavy-tails:rank}{corrélation de rang} plutôt que la corrélation de Pearson, et que ne vous apprennent ni l'une ni l'autre~?
\end{interviewq}

\begin{interviewq}[$\star\star\star$ \iqroles{chercheur, risque}]\label{iq:qm:robust-statistics-and-heavy-tails:6}
Que dit le théorème de Sklar, et pourquoi la \omterm{def:qm:robust-statistics-and-heavy-tails:copula}{copule} gaussienne sous-estime-t-elle les krachs conjoints~?
\end{interviewq}
